\section{Methods}\label{sec:methods}

\subsection{HEP implementation}
The HEP Registry is an append-only, event-sourced store.
Every operation appends an event that is stamped with its author and
hash-chained to the previous one, no event is overwritten or deleted,
and the current lifecycle state and belief of a hypothesis are derived
by replaying its event log, which is persisted in the run's session
directory.
Beyond the fields introduced in Results, each hypothesis records a
testable predicted observable and its lineage (parent hypotheses and
generation).
An \emph{inspired-by} link is recorded as provenance only
and does not make the new hypothesis a descendant.

The agent operates the registry through seven tools.
\texttt{propose\_hypothesis} registers a new hypothesis under a
research question, and \texttt{transition\_hypothesis} moves a
hypothesis between lifecycle states.
\texttt{attach\_evidence} records one piece of evidence on a
hypothesis and updates its belief from it.
\texttt{list\_hypotheses} returns the competing set grouped by
research question, with the current state and belief of each
hypothesis, and \texttt{get\_hypothesis} returns the full record of a
single hypothesis, including its event timeline.
Evolution is evidence-driven by construction:
\texttt{refine\_hypothesis} creates a child hypothesis while the
parent retains its own state, \texttt{merge\_hypotheses} combines
hypotheses on the same question and transitions unresolved parents to
\texttt{dormant} as subsumed, and both operations are accepted only
for parents that have reached a verdict or have at least one piece of
evidence attached, so that untested proposals cannot produce
descendants.

Beliefs are elicited rather than computed.
The agent states a prior
$P(H)$ with a written rationale at proposal, and at each evidence
attachment (with kind simulation, experiment, literature, derivation,
or analysis, and direction supports, refutes, or inconclusive) it
states an updated $P(H)$, recorded with its rationale and the
identifier of the triggering evidence.
The two belief rules described in Results are enforced by the tools
rather than left to the agent.
For computational and analytical evidence, the belief moves only if
the agent certifies the result as trustworthy for its method, with a
diagnostic covering both the adequacy of the setup and the validity of
the result.
An uncertified attachment is recorded but marked
insufficient.
Verdict transitions whose belief has not crossed the threshold, and
any move outside the lifecycle state machine, are rejected.

\subsection{Agent setup}
All experiments use a single autonomous research agent, implemented
with the OpenAI Agents SDK, that receives one task prompt and runs to
completion without human input.
Besides the HEP tools, the agent has computational tools for
atomistic materials research exposed as MCP servers, including
MLIP-based structure relaxation, molecular dynamics, and
materials-database queries, together with a library of documented
skills it can discover and execute, both provided by
AtomisticSkills~\cite{Deng2026AtomisticSkills}, as well as a tool for
authoring and executing its own analysis code, with each script
retained in the run's record, and recording tools for logging findings
and writing the final report.
The agent runs in an outer loop.
Each iteration allows up to 80 LLM
turns, after which the driver re-prompts it to continue, appending the
remaining wall-clock budget so that the agent can select tests that
fit within the remaining time.
The conversation history is persisted,
and every run is given a 24-hour wall-clock budget, which no run
exhausted.
The HEP-equipped agent is given the HEP tools and a
hypothesis-discipline prompt but no planning tools.
The run ends when
the investigation saturates, that is, when no tool call is made for
three consecutive iterations, at which point the agent is prompted to
consolidate its findings and write the final report.

\subsection{Trajectory analysis}
For the loop analysis in Fig.~\ref{fig:hep_loop}c, each step of a run
is classified into the loop elements.
HEP tool calls map deterministically: \texttt{propose\_hypothesis} to
H; \texttt{refine\_hypothesis} and \texttt{merge\_hypotheses} to U;
\texttt{attach\_evidence} to E, and additionally to B when the
attachment moved the belief; and \texttt{transition\_hypothesis} to J
for verdict transitions.
Calls to computational tools map to T, free-text steps are classified
by an LLM (GPT-5.5), and scaffolding steps such as file management are
excluded.
Consecutive repeats of the same element are compressed, and the
transition matrices count adjacent pairs of elements, row-normalized
and averaged over runs.
